% !TeX root = ../report_jouleecosystem.tex
\chapter{Conclusion}

This report documented the three interconnected parts of the energy-aware deep learning ecosystem built during this internship, following the data flow from physical silicon to training loops:

\textbf{JouleQuest} turned a fragile, manual measurement procedure into a robust, automated pipeline. Using single-command execution, adaptive burst sizing, and host-target clock synchronization, its campaign launcher profiles hundreds of layer configurations over multi-day acquisition sessions and resumes automatically after interruptions. It delivers empirical energy lookup tables for \texttt{Linear}, \texttt{Conv2d}, and \texttt{Attention} operators across diverse edge devices.

\textbf{JouleGrad} transforms these empirical tables into continuous, differentiable surrogates inside PyTorch training loops. Its multilinear interpolation produces analytical energy gradients without discrete noise, while strict out-of-range policies safeguard against unmeasured extrapolation artifacts.

\textbf{JouleNAS} guides neural architecture search with physical energy gradients, pruning redundant channels from deep models while balancing task accuracy against hardware power. Sweeping the energy penalty establishes empirical Pareto frontiers, while an inverse scaling law enables reliable hyperparameter transfer across different hardware targets.

Finally, \textbf{The Bridge} validates the discovered sub-networks back on physical silicon through a spec-driven materialization pipeline. Deploying real, structurally pruned models confirmed real-world energy reductions of up to $60.1\%$ on the Raspberry~Pi~5 and $20.0\%$ on the NVIDIA Jetson AGX Orin, comfortably preserving baseline accuracy.

The central takeaway of this work is that physical energy measurements expose hardware nuances that theoretical metrics cannot capture. With automated acquisition, differentiable estimation, hardware-guided search, and physical validation now unified end-to-end, the ecosystem provides a solid, accessible foundation for the next researcher to build upon.
